\section{Proyecto LINCE}

El Proyecto LINCE (Línea de industrialización de cargas de pago y plataformas espaciales) es
una iniciativa estratégica enmarcada en el PERTE Aeroespacial, cuyo propósito es consolidar
la capacidad industrial de España en el segmento de los microsatélites (100--200 kg)
destinados a órbitas bajas (LEO). Liderado por Indra, el consorcio busca desarrollar
plataformas satelitales modulares, fiables y de altas prestaciones, optimizando costes y
plazos mediante procesos de producción en serie. Este ecosistema integra a grandes
empresas, PYMES y centros de investigación, como la Universidad Politécnica de Madrid
(UPM), para habilitar servicios avanzados de telecomunicaciones, vigilancia y observación,
garantizando así la autonomía estratégica de España y Europa en el sector espacial
\cite{indra2026lince}.

\section{MPSoC Zynq UltraScale+ ZCU102}

La placa de desarrollo sobre la que se va a trabajar en este proyecto es la Tarjeta de
Evaluación de propósito general ZCU102, desarrollada por Xilinx, orientada al prototipado
rápido de sistemas empotrados de alta complejidad. Esta placa ofrece soporte a todas las
interfaces de alta velocidad y lógica programable del MPSoC (\textit{Multiprocessor System-on-Chip})
Zynq UltraScale+, que integra en un único encapsulado de silicio un sistema de
procesamiento con cuatro núcleos y una matriz lógica programable, conectados internamente
mediante buses de alto rendimiento. Esta arquitectura divide el sistema en dos dominios
complementarios, permitiendo separar la carga de trabajo entre el Sistema de Procesamiento
(PS, \textit{Processing System}) y la lógica programable (PL, \textit{Programmable Logic}).

El \textbf{PS} representa el dominio del software. Su desarrollo es análogo a la programación de
cualquier microcontrolador o microprocesador clásico; está basado en una arquitectura física
fija (núcleos ARM en este dispositivo) diseñada para ejecutar instrucciones de código en C o
C++ de forma estrictamente secuencial. Este dominio es el encargado de ejecutar sistemas
operativos (como RTEMS o Linux), gestionar la memoria y ejecutar las rutinas de control de
alto nivel.

Por el contrario, la \textbf{PL} representa el dominio del hardware a medida. Se trata de una matriz
tipo FPGA (\textit{Field Programmable Gate Array}), la cual no ejecuta un programa línea a línea,
sino que se reconfigura físicamente. Mediante lenguajes de descripción de hardware (como
VHDL), se define la interconexión de miles de recursos internos disponibles en el silicio
---tales como tablas de búsqueda (LUTs), biestables, memorias BRAM y bloques de
procesamiento digital de señales (DSP)--- para conformar circuitos digitales reales y
específicos. Esta característica otorga a la PL un determinismo temporal estricto y la
capacidad de procesar señales con un paralelismo masivo inalcanzable para un procesador
convencional.

Entre sus características destacadas, una que resulta relevante para este trabajo es la
capacidad de expansión externa mediante dos conectores FMC HPC (\textit{High Pin Count}). Estos
puertos son el punto de anclaje para tarjetas de expansión o \textit{Mezzanines}, permitiendo
extraer pines digitales y pares diferenciales hacia el exterior del sistema. Siguen el estándar
VITA 57.1 \textit{FPGA Mezzanine Card (FMC) specification} \cite{vita571fmc}.

\section{Vivado y Vitis (Xilinx): PS y PL}

Los programas que más se van a utilizar durante este trabajo son Vivado y Vitis, ambos
desarrollados por Xilinx para el desarrollo sobre sus sistemas hardware.

\subsection{Vivado Design Suite}

Es el entorno oficial de Xilinx para el diseño, síntesis e implementación de circuitos digitales
destinados a la PL. Permite programar en Lenguaje de Descripción Hardware (HDL) los
circuitos que se desean implementar, ofrece una herramienta de diseño por bloques y
permite realizar simulaciones temporales para verificar el comportamiento del hardware
desarrollado.

En el contexto específico de este proyecto, Vivado se utiliza para encapsular los transceptores
serie diseñados a medida en VHDL y enlazarlos al núcleo del procesador Zynq. Esta
integración se realiza mediante el bus de interconexión de alto rendimiento AXI
(\textit{Advanced eXtensible Interface}) a través de la herramienta \textit{Block Design}. Asimismo, Vivado
gestiona el archivo de restricciones físicas (.xdc), necesario para mapear lógicamente las
señales de los periféricos desde el silicio hacia los pines reales de los conectores FMC. El
flujo de trabajo en Vivado culmina con la generación del \textit{Bitstream} (el archivo binario que
configura físicamente la FPGA) y la exportación de la arquitectura del hardware en un
archivo unificado (.xsa).

\subsection{Vitis Unified Software Platform}

Vitis es el entorno de desarrollo integrado (IDE) proporcionado por Xilinx para la
programación del dominio del software, es decir, el PS. Funciona en tándem con Vivado:
importa el archivo de especificación hardware (.xsa) y genera de forma automatizada el
Paquete de Soporte de la Placa o BSP (\textit{Board Support Package}). Este BSP contiene el mapa
de memoria estático y los controladores de bajo nivel (\textit{drivers}) indispensables para que el
código en C/C++ pueda acceder a los periféricos instanciados previamente en la PL.

Durante el desarrollo de este trabajo, Vitis se emplea para la compilación cruzada orientada a
los núcleos ARM del sistema. Se utiliza para generar el firmware crítico de inicialización,
concretamente el PMUFW (\textit{Power Management Unit Firmware}) y el cargador de arranque de
primera etapa FSBL (\textit{First Stage Boot Loader}). Finalmente, Vitis proporciona las herramientas
de empaquetado para fusionar el código ejecutable, el \textit{Bitstream} de la PL y los archivos de
arranque en un único fichero binario (BOOT.BIN), habilitando el arranque autónomo de la
placa desde una memoria no volátil como una tarjeta SD.

\section{RTEMS}

RTEMS (\textit{Real-Time Executive for Multiprocessor Systems}) es un sistema operativo de tiempo
real (RTOS) de código abierto diseñado de forma específica para sistemas empotrados de
misión crítica. A diferencia de los sistemas operativos de propósito general, como
distribuciones estándar de Linux, RTEMS no utiliza memoria virtual y opera en un único
espacio de direcciones físico (arquitectura de un solo proceso y múltiples hilos). Esta
arquitectura elimina la latencia de los cambios de contexto pesados, minimiza la sobrecarga
del procesador y garantiza tiempos de respuesta estrictamente deterministas ante las
interrupciones del hardware.

El uso de RTEMS es un estándar de facto en la industria aeroespacial europea y americana
(siendo empleado activamente por agencias como la ESA y la NASA en misiones satelitales y
sondas espaciales) debido a su robustez, su certificación para entornos críticos y su
compatibilidad nativa con arquitecturas ARM como la del Zynq UltraScale+. En el marco del
proyecto LINCE y de este TFM, RTEMS se ejecuta sobre el PS actuando como el cerebro
temporal del sistema. Su función es orquestar las rutinas de control de alto nivel, gestionar
la lectura/escritura de los registros de los transceptores de la PL y asegurar que las
transacciones en los buses de comunicación (SpaceWire, RS422, RS485 o CAN) cumplan con
los plazos de tiempo (\textit{deadlines}) exigidos durante las pruebas \textit{Hardware-in-the-Loop} de los
subsistemas del satélite.

\section{Comunicaciones serie}

En el diseño de sistemas empotrados de misión crítica y arquitecturas satelitales, las
comunicaciones serie constituyen la columna vertebral para el intercambio de datos entre el
Ordenador de a Bordo (OBC) y sus múltiples periféricos, tales como sensores de actitud,
actuadores mecánicos y cargas de pago. Frente a las topologías de bus paralelo
tradicionales, los enlaces serie proporcionan una reducción drástica de las líneas físicas
necesarias, lo que se traduce de forma directa en una disminución crítica del volumen y peso
del arnés del satélite. Además, al operar frecuentemente mediante señalización diferencial,
mitigan en gran medida la susceptibilidad a las interferencias electromagnéticas (EMI)
propias del entorno espacial.

En este apartado se detallan los fundamentos técnicos, topologías y características eléctricas
de los protocolos de comunicación serie específicos que han sido seleccionados e integrados
en la plataforma de hardware durante el desarrollo de este trabajo: RS485, RS422, SpaceWire
y bus CAN.

\subsection{RS485}

El estándar TIA/EIA-485, comúnmente conocido como RS485, define las características
eléctricas de receptores y transmisores para su uso en sistemas de comunicaciones serie
balanceados (diferenciales). Su principal característica es la capacidad de operar en
topologías multipunto (buses compartidos), permitiendo conectar hasta 32 transceptores
estándar en el mismo par de cables (o más, si se emplean transceptores de carga
fraccional).

La señalización diferencial, donde la información se transmite como la diferencia de potencial
entre dos hilos trenzados (A y B), le otorga un alto Rechazo al Modo Común (CMRR). Esto
significa que cualquier ruido electromagnético inducido en el cable afecta por igual a ambas
líneas, cancelándose en el receptor. Habitualmente se implementa en modo
\textit{Half-Duplex} (2 hilos), donde la transmisión y recepción comparten el medio físico, exigiendo
un control estricto del flujo de datos mediante software (gestión de los pines de
\textit{Driver Enable}) para evitar colisiones. Es un estándar robusto capaz de alcanzar distancias
de hasta 1.200 metros (a velocidades reducidas) o velocidades de hasta 50 Mbps en
distancias cortas, requiriendo resistencias de terminación (típicamente de $120\,\Omega$) en
los extremos del bus para evitar reflexiones destructivas de la señal
\cite{ti2014rs485,eia485a}.

\subsection{RS422}

El estándar TIA/EIA-422 (RS422) es el predecesor técnico del RS485 y comparte con él la
base física de la señalización diferencial para garantizar la inmunidad al ruido y cubrir largas
distancias. Sin embargo, su principal diferencia radica en la topología de la red: el RS422
está diseñado estrictamente para comunicaciones punto a punto o punto a multipunto
(conocido como \textit{multi-drop}).

A diferencia del RS485, los transceptores RS422 no están diseñados para ceder el control del
bus. En un bus RS422 solo puede existir un único transmisor (maestro) conectado a un
máximo de 10 receptores (esclavos) en el mismo par de hilos. Para lograr una comunicación
bidireccional continua (\textit{Full-Duplex}), el RS422 requiere obligatoriamente el uso de cuatro
hilos (dos pares trenzados: uno dedicado exclusivamente a la transmisión y otro a la
recepción). Esta separación física de los canales de ida y vuelta elimina la necesidad de
arbitrar el bus por software, reduciendo la sobrecarga de procesamiento y garantizando un
determinismo temporal absoluto, lo que lo hace idóneo para el envío continuo de telemetría
de sensores críticos \cite{ti2018rs422,tia422b}.

\subsection{SpaceWire}

SpaceWire es un estándar de red de comunicaciones espaciales de alta velocidad, coordinado
por la Agencia Espacial Europea (ESA) y ampliamente adoptado a nivel mundial en misiones
científicas y comerciales. Está diseñado específicamente para interconectar nodos de alto
rendimiento a bordo de satélites, como memorias masivas, procesadores de instrumentación
y enlaces de telemetría, ofreciendo un gran ancho de banda y una fiabilidad extrema.

A nivel físico, SpaceWire utiliza señalización diferencial de bajo voltaje (LVDS) e implementa
una codificación única denominada \textit{Data-Strobe} (DS). En lugar de transmitir una señal de
reloj independiente que sufriría desvíos temporales (\textit{skew}) respecto a los datos a altas
frecuencias, la señalización DS envía los datos por un par diferencial y una señal
\textit{Strobe} por otro. El reloj se recupera en el receptor realizando una operación lógica XOR
entre la señal de datos y el \textit{Strobe}. Esto permite enlaces punto a punto \textit{full-duplex}
asíncronos con velocidades que abarcan desde los 2~Mbps hasta más de 400~Mbps. Además,
su arquitectura en capas soporta el enrutamiento de paquetes, permitiendo construir
topologías de red complejas mediante \textit{routers} SpaceWire \cite{ecss50spw,parkes2012spw}.

\subsection{CAN}

El bus CAN es un estándar de comunicaciones serie robusto diseñado originalmente para el
sector de la automoción, pero cuya alta tolerancia a fallos lo ha convertido en un estándar
de facto (\textit{CAN for Aerospace}) para las redes internas de monitorización, telemetría y
telecomandos (TM/TC) en nanosatélites y satélites comerciales.

Es un bus multipunto y multi-maestro que opera bajo el paradigma CSMA/CD+AMP
(\textit{Carrier Sense Multiple Access with Collision Detection and Arbitration on Message
Priority}). En una red CAN, los nodos no tienen una dirección física; en su lugar, transmiten
tramas de datos encabezadas por un identificador que define el contenido y la prioridad del
mensaje. Si dos nodos intentan transmitir simultáneamente, la colisión se resuelve a nivel de
bit y de forma no destructiva: el mensaje con el identificador de mayor prioridad sobrescribe
al de menor prioridad sin corromper el bus, garantizando un determinismo estricto para las
alarmas críticas. Además, la capa de enlace del protocolo CAN incorpora mecanismos de
detección de errores por hardware (CRC, \textit{bit stuffing}, comprobación de tramas) y
confinamiento automático de nodos defectuosos, asegurando que un periférico averiado no
bloquee el sistema de control de actitud u otros subsistemas vitales del Ordenador de a Bordo
\cite{iso118981,esa2015can}.
